\chapter{The sheaf-of-modules over-equivalence (shared slice root)}
\label{chap:Picard_SheafOverEquivalence}
% archon:covers AlgebraicJacobian/Picard/SheafOverEquivalence.lean

\section{The equivalence}
\label{sec:soe_equivalence}

\begin{definition}[Sheaf-of-modules over-equivalence]\leanok
  \label{def:sheafofmodules_over_equivalence}
  \dcref{custom}
  \lean{AlgebraicGeometry.Scheme.Modules.overEquivalence}
  \uses{def:phi_over, def:psi_over, def:over_equiv_functor_is_continuous,
    def:over_equiv_inverse_is_continuous}
  % SOURCE: Mathlib code-location provenance (this is an Archon-original assembly
  % of existing Mathlib primitives, NOT an external literature citation, so it
  % carries no verbatim mathematical SOURCE QUOTE). The primitive instantiated is
  % SheafOfModules.pushforwardPushforwardEquivalence, defined at
  % Mathlib/Algebra/Category/ModuleCat/Sheaf/PushforwardContinuous.lean:305
  % (section Equivalence, lines 289--329): "equivalence of SheafOfModules
  % categories from a site equivalence e : C \simeq D together with a morphism of
  % the two structure ring sheaves and its inverse, with two coherences". The
  % underlying site equivalence is TopologicalSpace.Opens.overEquivalence
  % (Mathlib/Topology/Sheaves/Over.lean:41), and the continuity prerequisites are
  % the documented TODO at Over.lean:19--22.
  Let \(X\) be a scheme and \(U \subseteq X\) an open subset, with associated open
  immersion \(\iota : \widetilde{U} \hookrightarrow X\). The reindexing site
  equivalence
  \[
    e \;:=\; \mathtt{Opens.overEquivalence}\,U
       \;:\; \mathtt{Over}\,U \;\simeq\; \mathtt{Opens}\,\widetilde{U}
  \]
  --- between the slice of \(\mathtt{Opens}\,X\) over \(U\) (carrying the slice
  topology \((\mathtt{Opens.grothendieckTopology}\,X).\mathtt{over}\,U\)) and the
  opens of the subspace \(\widetilde{U}\) (carrying
  \(\mathtt{Opens.grothendieckTopology}\,\widetilde{U}\)) --- lifts to an
  \emph{equivalence of sheaf-of-modules categories}
  \[
    \mathtt{overEquivalence}\;:\;
      \mathtt{SheafOfModules}\,\bigl((\widetilde{U}).\mathtt{ringCatSheaf}\bigr)
        \;\;\simeq\;\;
      \mathtt{SheafOfModules}\,\bigl(X.\mathtt{ringCatSheaf}.\mathtt{over}\,U\bigr),
  \]
  the modules-level lift of \(e\) that transports the structure ring sheaf. It is
  obtained by instantiating
  \(\mathtt{SheafOfModules.pushforwardPushforwardEquivalence}\) at \(e\), with the
  open-immersion structure-sheaf ring morphism
  \[
    \varphi \;:\;
      X.\mathtt{ringCatSheaf}.\mathtt{over}\,U
      \;\longrightarrow\;
      \bigl(e.\mathtt{functor}.\mathtt{sheafPushforwardContinuous}\bigr).\mathrm{obj}\,
        \bigl((\widetilde{U}).\mathtt{ringCatSheaf}\bigr)
  \]
  --- sectionwise at an object \(V \in \mathtt{Over}\,U\) the open-immersion
  comparison \(\mathcal{O}_X(V.\mathrm{left}) \to \mathcal{O}_{\widetilde{U}}(e\,V)\),
  i.e.\ the section-level isomorphism \(\iota.\mathtt{appIso}\) of the structure
  sheaves --- together with its inverse \(\psi\) and the two coherences
  \(H_1, H_2\) expressing that \(\varphi\) and \(\psi\) are mutually inverse.
\end{definition}

\begin{proof}

\leanok
  \uses{def:over_equiv_functor_is_continuous, def:over_equiv_inverse_is_continuous,
        def:phi_over, def:psi_over}
  The construction is the assembly of three existing Mathlib primitives; the only
  step carrying mathematical content is the ring morphism \(\varphi\).

  \emph{Continuity (no work).} The equivalence
  \(\mathtt{pushforwardPushforwardEquivalence}\) requires both legs of the site
  equivalence \(e\) to be continuous:
  \(e.\mathtt{functor}\) and \(e.\mathtt{inverse}\). For a dense-subsite inclusion
  continuity is automatic --- it is a typeclass instance --- and for an
  equivalence the density of one leg propagates to the other. The declaration
  \(\mathtt{overEquivInverseIsDenseSubsite}\) exhibits \(e.\mathtt{inverse}\) as a
  dense subsite of the slice site, from which both continuity legs follow by
  inference. The Mathlib TODO at \texttt{Topology/Sheaves/Over.lean} (lines 19--22)
  is discharged for this case.

  \emph{The ring morphism \(\varphi\) (the genuine content).} The two structure
  ring sheaves are different objects on different sites:
  \(X.\mathtt{ringCatSheaf}.\mathtt{over}\,U\) over \(\mathtt{Over}\,U\), and
  \((\widetilde{U}).\mathtt{ringCatSheaf}\) over \(\mathtt{Opens}\,\widetilde{U}\).
  The open immersion \(\iota\) identifies them: at an object \(V \in
  \mathtt{Over}\,U\) the sections \(\mathcal{O}_X(V.\mathrm{left})\) and
  \(\mathcal{O}_{\widetilde{U}}(\iota^{-1}(V.\mathrm{left}))\) agree through the
  open-immersion section isomorphism \(\iota.\mathtt{appIso}\). This is verbatim the
  datum that \(\mathtt{Scheme.Modules.restrictFunctor}\) already uses inline
  (its structure map is \((\iota.\mathtt{appIso}\,\cdot).\mathrm{inv}\)). Packaging
  this section-level isomorphism naturally in \(V\) yields \(\varphi\); its inverse
  is \(\psi\), and the coherences \(H_1, H_2\) record that the
  \(\iota.\mathtt{appIso}\) round-trips are identities, proved by the
  inverse-iso laws (\(\mathtt{Sheaf.hom\_ext}\) reduces them sectionwise).

  Feeding \(e\), \(\varphi\), \(\psi\), \(H_1\), \(H_2\) into
  \(\mathtt{pushforwardPushforwardEquivalence}\) produces the asserted equivalence.
  Its underlying functor is \(\mathtt{SheafOfModules.pushforward}\,\varphi\), which
  is what the two consumer isomorphisms below unfold.
\end{proof}

\subsection{Substrate helpers for the equivalence}
\label{sec:soe_equivalence_helpers}

\begin{definition}[Continuity of the inverse leg of the over-equivalence]\leanok
  \label{def:over_equiv_inverse_is_continuous}
  \dcref{custom}
  \lean{AlgebraicGeometry.Scheme.Modules.overEquivInverseIsContinuous}
  % NOTE iter-283: removed the spurious `\uses{def:sheafofmodules_over_equivalence}`.
  % This continuity fact is a PREREQUISITE that `pushforwardPushforwardEquivalence`
  % demands in order to BUILD the over-equivalence -- it is about the site
  % equivalence `Opens.overEquivalence U`'s inverse functor (Lean body:
  % `(Opens.overEquivalence U).inverse.IsContinuous ...`, closed by
  % `infer_instance`) and does not reference the sheaf-of-modules `overEquivalence`.
  % Listing the equivalence here reversed the dependency and created an
  % overEquivalence <-> inverse_is_continuous cycle that crashed `leanblueprint web`.
  % It has no blueprint-level dependency (the site equivalence and Grothendieck
  % topologies are Mathlib primitives), so the \uses annotation is dropped entirely
  % rather than left empty (an empty \uses{} is itself a malformed annotation).
  The inverse leg \(e.\mathtt{inverse} : \mathtt{Opens}\,\widetilde{U} \to
  \mathtt{Over}\,U\) of the reindexing site equivalence
  \(e = \mathtt{Opens.overEquivalence}\,U\) is continuous from the subspace
  topology \(\mathtt{Opens.grothendieckTopology}\,\widetilde{U}\) to the slice
  topology \((\mathtt{Opens.grothendieckTopology}\,X).\mathtt{over}\,U\). This is the
  continuity prerequisite, recorded for the scheme-carrier presentation of the
  open subscheme, that \(\mathtt{pushforwardPushforwardEquivalence}\) demands of the
  inverse leg.
\end{definition}

\begin{proof}

\leanok
  Proved directly in Lean: it converts to the bare-subspace form
  \(\widetilde{U}\), where the project's dense-subsite instance for the inverse leg
  and the priority instance \(\mathtt{IsDenseSubsite} \Rightarrow
  \mathtt{IsContinuous}\) discharge it by inference.
\end{proof}

\begin{definition}[Continuity of the functor leg of the over-equivalence]\leanok
  \label{def:over_equiv_functor_is_continuous}
  \dcref{custom}
  \lean{AlgebraicGeometry.Scheme.Modules.overEquivFunctorIsContinuous}
  \uses{def:over_equiv_inverse_is_continuous}
  The functor leg \(e.\mathtt{functor} : \mathtt{Over}\,U \to
  \mathtt{Opens}\,\widetilde{U}\) of the reindexing site equivalence is continuous
  from the slice topology \((\mathtt{Opens.grothendieckTopology}\,X).\mathtt{over}\,U\)
  to the subspace topology \(\mathtt{Opens.grothendieckTopology}\,\widetilde{U}\). It
  is the second continuity prerequisite of
  \(\mathtt{pushforwardPushforwardEquivalence}\).
\end{definition}

\begin{proof}
  \uses{def:over_equiv_inverse_is_continuous}
  \leanok
  Proved directly in Lean. Converting to the bare-subspace form, the density of the
  inverse leg (\cref{def:over_equiv_inverse_is_continuous}) propagates, for an
  equivalence of sites, to density --- hence continuity --- of the functor leg.
\end{proof}

\begin{lemma}[Reindexed open recovers its left component]\leanok
  \label{lem:image_overequiv_functor_obj}
  \dcref{custom}
  \lean{AlgebraicGeometry.Scheme.Modules.image_overEquiv_functor_obj}
  For an object \(V \in \mathtt{Over}\,U\), the open-immersion image of the
  reindexed open \(e.\mathtt{functor}.\mathrm{obj}\,V = \iota^{-1}(V.\mathrm{left})\)
  is \(V.\mathrm{left}\) itself,
  \[
    \iota_{\ast}\,\bigl(e.\mathtt{functor}.\mathrm{obj}\,V\bigr)
      \;=\; V.\mathrm{left},
  \]
  because \(V.\mathrm{left} \le U\). This is the equality of opens that identifies
  the two structure-sheaf presheaves underlying the ring morphism \(\varphi\).
\end{lemma}

\begin{proof}

\leanok
  Proved directly in Lean: a point \(y\) lies in
  \(\iota_{\ast}(\iota^{-1}(V.\mathrm{left}))\) iff it lies in
  \(V.\mathrm{left}\), using \(V.\mathrm{left} \le U\) so that the preimage and
  image cancel.
\end{proof}

\begin{lemma}[Inverse-reindexed open is its open-immersion image]\leanok
  \label{lem:left_overequiv_inverse_obj}
  \dcref{custom}
  \lean{AlgebraicGeometry.Scheme.Modules.left_overEquiv_inverse_obj}
  For an open \(W \subseteq \widetilde{U}\), the left component of the
  inverse-reindexed slice object \(e.\mathtt{inverse}.\mathrm{obj}\,W\) is the
  open-immersion image of \(W\),
  \[
    \bigl(e.\mathtt{inverse}.\mathrm{obj}\,W\bigr).\mathrm{left}
      \;=\; \iota_{\ast}\,W .
  \]
  This is the symmetric counterpart of \cref{lem:image_overequiv_functor_obj}.
\end{lemma}

\begin{proof}

\leanok
  Proved directly in Lean: the two opens have equal underlying sets by definitional
  unfolding of the inverse reindexing functor.
\end{proof}

\begin{definition}[The structure-sheaf ring morphism \(\varphi\)]\leanok
  \label{def:phi_over}
  \dcref{custom}
  \lean{AlgebraicGeometry.Scheme.Modules.phiOver}
  \uses{lem:image_overequiv_functor_obj}
  The ring morphism \(\varphi\) underlying the over-equivalence: at an object
  \(V \in \mathtt{Over}\,U\) it is the identification
  \(\mathcal{O}_X(V.\mathrm{left}) \xrightarrow{\sim}
  \mathcal{O}_{\widetilde{U}}(\iota^{-1}(V.\mathrm{left}))\) of the two structure
  sheaves, realised on the nose as \(X.\mathtt{ringCatSheaf}\) applied to the open
  equality \cref{lem:image_overequiv_functor_obj} (the two presheaves agreeing
  definitionally). Naturality in \(V\) follows from functoriality of
  \(X.\mathtt{ringCatSheaf}\).
\end{definition}

\begin{proof}
  \uses{lem:image_overequiv_functor_obj}
  \leanok
  Proved directly in Lean. It is a sub-step of
  \cref{def:sheafofmodules_over_equivalence}.
\end{proof}

\begin{definition}[The inverse structure-sheaf ring morphism \(\psi\)]\leanok
  \label{def:psi_over}
  \dcref{custom}
  \lean{AlgebraicGeometry.Scheme.Modules.psiOver}
  \uses{lem:left_overequiv_inverse_obj}
  The inverse ring morphism \(\psi\), symmetric to \cref{def:phi_over}: at an open
  \(W \subseteq \widetilde{U}\) it is \(X.\mathtt{ringCatSheaf}\) applied to the open
  equality \cref{lem:left_overequiv_inverse_obj}, providing the section-level inverse
  of \(\varphi\). It is the datum \(\psi\) fed into
  \(\mathtt{pushforwardPushforwardEquivalence}\).
\end{definition}

\begin{proof}
  \uses{lem:left_overequiv_inverse_obj}
  \leanok
  Proved directly in Lean. It is a sub-step of
  \cref{def:sheafofmodules_over_equivalence}.
\end{proof}

\section{The equivalence carries restriction and unit}
\label{sec:soe_consumers}

\subsection{Substrate helpers for the restriction comparison}
\label{sec:soe_consumers_helpers}

\begin{definition}[The restriction ring morphism along the open-immersion functor]\leanok
  \label{def:psi_restrict}
  \dcref{custom}
  \lean{AlgebraicGeometry.Scheme.Modules.psiRestrict}
  The ring morphism along the open-immersion functor
  \(\iota_{\mathtt{opensFunctor}} : \mathtt{Opens}\,\widetilde{U} \to
  \mathtt{Opens}\,X\) underlying \(\Scheme.\mathtt{Modules.restrictFunctor}\,\iota\):
  at an open \(W \subseteq \widetilde{U}\) it is the inverse open-immersion section
  comparison \((\iota.\mathtt{appIso}\,W)^{-1}\), promoted to a ring-sheaf morphism.
  It is reconstructed so that the identification
  \cref{lem:restrict_functor_eq_pushforward_psi_restrict} holds definitionally.
\end{definition}

\begin{proof}

\leanok
  Proved directly in Lean: it repackages the structure map of
  \(\mathtt{restrictFunctor}\,\iota\) (whose component at \(W\) is
  \((\iota.\mathtt{appIso}\,W)^{-1}\)) as a morphism of ring sheaves.
\end{proof}

\begin{lemma}[Restriction functor as a pushforward]\leanok
  \label{lem:restrict_functor_eq_pushforward_psi_restrict}
  \dcref{custom}
  \lean{AlgebraicGeometry.Scheme.Modules.restrictFunctor_eq_pushforward_psiRestrict}
  \uses{def:psi_restrict}
  The restriction functor along the open immersion is, on the nose, a sheaf-of-modules
  pushforward along the morphism \cref{def:psi_restrict}:
  \[
    \Scheme.\mathtt{Modules.restrictFunctor}\,\iota
      \;=\; \mathtt{SheafOfModules.pushforward}\,(\mathtt{psiRestrict}).
  \]
\end{lemma}

\begin{proof}
  \uses{def:psi_restrict}
  \leanok
  Proved directly in Lean: the two functors are definitionally equal, since
  \(\mathtt{psiRestrict}\) was built to reproduce the internal structure map of
  \(\mathtt{restrictFunctor}\,\iota\) verbatim.
\end{proof}

\begin{definition}[Index reconciliation natural isomorphism]\leanok
  \label{def:over_forget_nat_iso}
  \dcref{custom}
  \lean{AlgebraicGeometry.Scheme.Modules.overForgetNatIso}
  \uses{lem:image_overequiv_functor_obj}
  The natural isomorphism
  \[
    \mathtt{Over.forget}\,U \;\xrightarrow{\ \sim\ }\;
      e.\mathtt{functor} \;\circ\; \iota_{\mathtt{opensFunctor}}
  \]
  reconciling the two index functors \(\mathtt{Over}\,U \to \mathtt{Opens}\,X\)
  entering the composite pushforward; both send \(V \mapsto V.\mathrm{left}\), and the
  component at \(V\) is the \(\mathtt{eqToIso}\) of the open equality
  \cref{lem:image_overequiv_functor_obj}.
\end{definition}

\begin{proof}
  \uses{lem:image_overequiv_functor_obj}
  \leanok
  Proved directly in Lean: the components are \(\mathtt{eqToIso}\) of
  \cref{lem:image_overequiv_functor_obj}, and naturality is automatic since the opens
  form a thin category.
\end{proof}

\begin{lemma}[The over-equivalence carries restriction to over]\leanok
  \label{lem:sheafofmodules_restrict_over_iso}
  \dcref{custom}
  \lean{AlgebraicGeometry.Scheme.Modules.restrictOverIso}
  \uses{def:sheafofmodules_over_equivalence, def:phi_over, def:psi_restrict,
    def:over_forget_nat_iso}
  Let \(M \in \Scheme.\mathtt{Modules}\,X\). Under the over-equivalence of
  \cref{def:sheafofmodules_over_equivalence} there is a canonical isomorphism
  \[
    \mathtt{overEquivalence}.\mathtt{functor}.\mathrm{obj}\,\bigl(M|_\iota\bigr)
      \;\xrightarrow{\ \sim\ }\;
    M.\mathtt{over}\,U,
  \]
  i.e.\ the equivalence carries the restriction \(M|_\iota\) of \(M\) along the
  open immersion to the slice object \(M.\mathtt{over}\,U\).
\end{lemma}

\begin{proof}
  \uses{def:sheafofmodules_over_equivalence, def:phi_over, def:psi_restrict,
        lem:restrict_functor_eq_pushforward_psi_restrict, def:over_forget_nat_iso,
        def:over_equiv_functor_is_continuous}
  \leanok
  The restriction \(M|_\iota\) is \emph{itself} a pushforward of sheaves of
  modules: by construction \(M|_\iota = \mathtt{SheafOfModules.pushforward}\)
  applied along \(\iota.\mathtt{opensFunctor}\) (the open-immersion functor
  \(\mathtt{Opens}\,\widetilde{U} \to \mathtt{Opens}\,X\)). The functor underlying
  the over-equivalence is \(\mathtt{SheafOfModules.pushforward}\,\varphi\). Their
  composite is again a single pushforward, identified by
  \(\mathtt{SheafOfModules.pushforwardComp}\), which here is the identity
  isomorphism \(\mathtt{Iso.refl}\). It remains to identify the two underlying
  index functors \(\mathtt{Over}\,U \to \mathtt{Opens}\,X\) entering the two
  pushforwards: both send an object \(V\) to its left component \(V.\mathrm{left}\),
  so they are equal, and \(\mathtt{pushforwardNatIso}\) along the \(\mathtt{eqToIso}\)
  of that equality transports the composite to \(M.\mathtt{over}\,U\). The whole
  argument mirrors, step for step, the existing
  \(\mathtt{Scheme.Modules.restrictFunctorAdjCounitIso}\).
\end{proof}

\begin{lemma}[The over-equivalence carries unit to unit]\leanok
  \label{lem:sheafofmodules_unit_over_iso}
  \dcref{custom}
  \lean{AlgebraicGeometry.Scheme.Modules.unitOverIso}
  \uses{def:sheafofmodules_over_equivalence, def:phi_over}
  Under the over-equivalence of \cref{def:sheafofmodules_over_equivalence} there is
  a canonical isomorphism
  \[
    \mathtt{overEquivalence}.\mathtt{functor}.\mathrm{obj}\,
      \bigl(\mathtt{SheafOfModules.unit}\,(\widetilde{U}).\mathtt{ringCatSheaf}\bigr)
      \;\xrightarrow{\ \sim\ }\;
    \mathtt{SheafOfModules.unit}\,\bigl(X.\mathtt{ringCatSheaf}.\mathtt{over}\,U\bigr),
  \]
  i.e.\ the equivalence carries the unit module (the structure sheaf-as-module) of
  the subspace site to the unit module of the slice site.
\end{lemma}

\begin{proof}
  \uses{def:sheafofmodules_over_equivalence}
  \leanok
  The functor underlying the over-equivalence is
  \(\mathtt{SheafOfModules.pushforward}\,\varphi\). The pushforward of a unit module
  along a morphism of ringed sites is the unit module of the target ring sheaf,
  transported along the ring identification \(\varphi\): the structure sheaf
  pushes forward to the structure sheaf because \(\varphi\) is precisely the
  open-immersion identification of the two structure ring sheaves. Concretely the
  pushforward-of-unit computation produces \(\mathtt{unit}\) of the codomain ring
  sheaf up to \(\varphi\), and \(\varphi\) being an isomorphism makes this a
  genuine isomorphism.

  Concretely, the canonical comparison map realising this transport is the
  unit-to-pushforward-of-unit morphism
  \[
    \mathtt{SheafOfModules.unitToPushforwardObjUnit}\,\varphi
      \;:\;
      \mathtt{unit}\,\bigl(X.\mathtt{ringCatSheaf}.\mathtt{over}\,U\bigr)
      \;\longrightarrow\;
      (\mathtt{pushforward}\,\varphi).\mathrm{obj}\,
        \bigl(\mathtt{unit}\,(\widetilde{U}).\mathtt{ringCatSheaf}\bigr).
  \]
  Its section at an open \(W\) is exactly the section-level ring map
  \(\varphi.\mathrm{app}\,W\): the morphism acts on a section by multiplication
  through \(\varphi\), so on \(W\) it is the additive-group map underlying
  \(\varphi.\mathrm{app}\,W\) (the characterisation
  \(\mathtt{unitToPushforwardObjUnit\_val\_app\_apply}\)). Since \(\varphi\) is an
  isomorphism of ring sheaves, each component \(\varphi.\mathrm{app}\,W\) is a ring
  isomorphism, hence its underlying additive-group map is an isomorphism.

  It follows that \(\mathtt{unitToPushforwardObjUnit}\,\varphi\) is an isomorphism.
  One reflects the sectionwise isomorphism upward through the two forgetful
  functors --- first \(\mathtt{SheafOfModules.forget}\) (down to presheaves of
  modules), then \(\mathtt{PresheafOfModules.toPresheaf}\) (down to presheaves of
  additive groups) --- so that being an isomorphism reduces to the componentwise
  isomorphism check
  \(\mathtt{NatTrans.isIso\_iff\_isIso\_app}\) on the underlying natural
  transformation, which is the sectionwise statement just established. The desired
  isomorphism \(\mathtt{unitOverIso}\) is then the inverse of the resulting
  isomorphism, \((\mathtt{asIso}\,(\mathtt{unitToPushforwardObjUnit}\,\varphi))^{-1}\).
  This is the slice-site analogue of
  \(\mathtt{pullbackObjUnitToUnitIso}\).
\end{proof}

\section{The engine corollary}
\label{sec:soe_chart}

\begin{lemma}[Over--restrict trivialisation bridge]\leanok
  \label{lem:chart_over_iso}
  \dcref{custom}
  \lean{AlgebraicGeometry.Scheme.Modules.chartOverIso}
  \uses{lem:sheafofmodules_restrict_over_iso, lem:sheafofmodules_unit_over_iso}
  Let \(M \in \Scheme.\mathtt{Modules}\,X\) and \(U \subseteq X\) open, and suppose
  given a trivialisation of the restriction along the open immersion
  \(\iota : \widetilde{U} \hookrightarrow X\),
  \[
    e \;:\; M|_\iota \;\xrightarrow{\ \sim\ }\;
      \mathtt{SheafOfModules.unit}\,(\widetilde{U}).\mathtt{ringCatSheaf}.
  \]
  Then there is a canonical isomorphism in the slice category
  \[
    \mathtt{chartOverIso}\,M\,U\,e \;:\;
      M.\mathtt{over}\,U \;\xrightarrow{\ \sim\ }\;
      \mathtt{SheafOfModules.unit}\,\bigl(X.\mathtt{ringCatSheaf}.\mathtt{over}\,U\bigr).
  \]
\end{lemma}

\begin{proof}
  \uses{lem:sheafofmodules_restrict_over_iso, lem:sheafofmodules_unit_over_iso,
        def:sheafofmodules_over_equivalence}
        \leanok
  Transport the trivialisation \(e\) across the over-equivalence and reconcile the
  two ends with the consumer isomorphisms. Explicitly, \(\mathtt{chartOverIso}\,M\,U\,e\)
  is the composite
  \[
    M.\mathtt{over}\,U
      \xrightarrow[\sim]{\ (\mathtt{restrictOverIso})^{-1}\ }
      \mathtt{overEquivalence}.\mathtt{functor}.\mathrm{obj}\,\bigl(M|_\iota\bigr)
      \xrightarrow[\sim]{\ \mathtt{overEquivalence}.\mathtt{functor}.\mathtt{mapIso}\,e\ }
      \mathtt{overEquivalence}.\mathtt{functor}.\mathrm{obj}\,(\mathtt{unit}\,\cdots)
      \xrightarrow[\sim]{\ \mathtt{unitOverIso}\ }
      \mathtt{unit}\,\bigl(X.\mathtt{ringCatSheaf}.\mathtt{over}\,U\bigr).
  \]
  The first arrow is the inverse of \cref{lem:sheafofmodules_restrict_over_iso}; the
  middle arrow is the image of \(e\) under the (functor of the) over-equivalence,
  which is functorial in isomorphisms hence sends the iso \(e\) to an iso; the last
  arrow is \cref{lem:sheafofmodules_unit_over_iso}. Each factor is an isomorphism,
  so the composite is.

  This lemma addresses the site mismatch in the chart-presentation of
  \cref{lem:lbc_chart_presentation}: the given trivialisation \(e\) lives in
  \(\mathtt{SheafOfModules}\,(\widetilde{U}).\mathtt{ringCatSheaf}\) (the open
  subspace site), whereas the finite-presentation transport
  \(\mathtt{Presentation.ofIsIso}\) must act in
  \(\mathtt{SheafOfModules}\,(X.\mathtt{ringCatSheaf}.\mathtt{over}\,U)\) (the slice
  site); \(\mathtt{chartOverIso}\) crosses that gap. It is stated in full
  generality in \(M\), \(U\), \(e\) for use in
  \cref{chap:Picard_LineBundleCoherence}. The over-equivalence of
  \cref{def:sheafofmodules_over_equivalence} likewise underlies the internal-Hom
  theory: the commutativity of the internal Hom with the slice-site change of base
  (\cref{lem:dual_restrict_iso}) passes through the same equivalence.
\end{proof}

\begin{definition}[Engine bridge: line-bundle over--restrict trivialisation]\leanok
  \label{def:linebundle_chart_over_iso}
  \dcref{custom}
  \lean{AlgebraicGeometry.Scheme.LineBundle.chartOverIso}
  \uses{lem:chart_over_iso}
  The line-bundle engine's local over--restrict trivialisation bridge: given
  \(M \in \Scheme.\mathtt{Modules}\,X\), an open \(U \subseteq X\), and a
  trivialisation \(e : M|_\iota \xrightarrow{\sim}
  \mathtt{SheafOfModules.unit}\,(\widetilde{U}).\mathtt{ringCatSheaf}\), it produces
  the slice-level isomorphism \(M.\mathtt{over}\,U \xrightarrow{\sim}
  \mathtt{SheafOfModules.unit}\,(X.\mathtt{ringCatSheaf}.\mathtt{over}\,U)\). It is the
  consumer-facing alias used by the finite-presentation engine, defined to be
  precisely the general construction \cref{lem:chart_over_iso}.
\end{definition}

\begin{proof}
  \uses{lem:chart_over_iso}
  \leanok
  Proved directly in Lean: the declaration is definitionally the general
  \(\Scheme.\mathtt{Modules.chartOverIso}\) of \cref{lem:chart_over_iso} applied to
  the same data \(M\), \(U\), \(e\).
\end{proof}
